\section{Introduction}

An index is a redundant structure: everything it contains is derivable from the
table it indexes. That redundancy is what makes it fast, and it is also what
makes its failures quiet. If the index and the table disagree, the database does
not stop; it answers a query with a subset of the correct answer, and no layer of
the system is in a position to notice.

This is not a hypothetical failure mode. Three defects from our own work on
PostgreSQL's generalized search tree fall squarely into it. A comparator used in
sorted construction was inconsistent for keys that mapped to equal values,
producing a tree whose ordering assumptions did not hold. Deleted pages could be
returned to circulation before every scan that might reach them had finished, so
a scan could read a page holding unrelated data. And building an index
concurrently did not account for prepared transactions, so rows committed through
two-phase commit could be silently absent from it. In each case every individual
operation returned a correct result. What was violated was a property of the
structure as a whole, and no test of individual operations can observe such a
property.

The conclusion we draw is that a class of correctness properties in a storage
engine needs a different discipline from testing behaviour: state the invariants
of the structure, and check them against the actual bytes on disk. That
conclusion is not novel in itself~--- it is the premise of amcheck for the
B-tree~\cite{thread-amcheck-btree}. What this paper contributes is its extension
to an extensible index framework, where the key type is supplied by the user and
the checker cannot know what a key means, and its extension to production
conditions, where the check must run without stopping the workload.

\paragraph{Contributions.}

\begin{enumerate}
\item Structural invariants for the generalized search tree and the generalized
      inverted index, expressed entirely through the operator class interface and
      therefore valid for any key type (Section~\ref{sec:invariants}).
\item A level-order traversal that verifies parent key consistency with memory
      bounded by the width of one tree level rather than by the size of the index
      (Section~\ref{sec:parent}).
\item Sibling-link verification with lock coupling, and the argument that makes
      coupling admissible in a checker (Section~\ref{sec:siblings}).
\item Normalization of index tuples, without which comparison against a
      regenerated tuple compares representations rather than values
      (Section~\ref{sec:normalize}).
\item An account of running the check on a physical replica, and of the defects
      found in the checker itself after release (Sections~\ref{sec:replica}
      and~\ref{sec:checker}).
\end{enumerate}

\section{Why the operator class interface is the right vocabulary}
\label{sec:invariants}

A generalized search tree~\cite{hellerstein1995generalized} is a balanced tree
whose node occupies a page and whose entries pair a key with a downlink, the key
covering every key in the subtree below. The key type and its operations belong
to an operator class supplied outside the framework: \code{consistent},
\code{union}, \code{penalty}, \code{picksplit} and the rest. The framework does
not know whether a key is a rectangle, a signature or a range, and a verification
tool built into the framework cannot know either.

This is a constraint, and it is also the design principle. If the invariants are
stated using only the operator class interface, the checker inherits the
extensibility of the framework: it works for a key type written after the checker
was. If they are stated in terms of any particular key semantics, it does not.
Every invariant below is therefore expressed through \code{consistent} alone,
which is the operation the search itself uses.

\begin{enumerate}
\item \emph{Parent key consistency.} For every entry of an internal page, the
      key of that entry covers the keys of all entries on the page it points to:
      \code{consistent}$(k_{\mathrm{parent}}, k_{\mathrm{child}})$ holds for
      every child. A violation means a search will not descend into a subtree
      that contains an answer~--- silently incomplete results, the failure mode
      of Section~1.
\item \emph{Balance.} All leaf pages lie at the same depth, and an internal page
      points either only to internal pages or only to leaf pages.
\item \emph{Completeness of the right-link chain.} A page's right link points to
      a page of the same level, and following the chain visits every page of the
      level exactly once.
\item \emph{Sibling agreement.} The right link of a left neighbour points to a
      page whose left link points back to it.
\item \emph{Unreachability of deleted pages.} No entry of any internal page
      points to a page marked deleted.
\end{enumerate}

Invariant~3 has a consequence for construction that is worth stating explicitly,
because it runs the other way from how such things are usually decided. Bulk
loading a tree bottom-up could leave right links unset, since no search needs
them: a search descends. Leaving them unset would force the checker to
distinguish indexes by how they were built, and an invariant that holds only for
some instances of a structure is not an invariant. The construction was therefore
made to set them. Verifiability is here a design input rather than an
afterthought.

For the inverted index the object is not a tree but a graph: an entry tree whose
leaves hold either inline posting lists or the roots of posting trees. The
invariants restate accordingly~--- key consistency along paths of the graph, at
least one downlink per internal page, uniformity of the level an internal page
points to, and ordering of entries within a page.

\section{Verifying parent keys without holding the index}
\label{sec:parent}

The direct way to check invariant~1 is to descend from the root, carrying each
parent key to its children. This is correct and unusable in production for two
reasons: the descent reads pages in an order unrelated to their placement, and it
holds state proportional to the width of the level being descended into, which
for a large index means holding most of the index.

The procedure used instead traverses in level order and passes parent keys
downward through intermediate storage. Memory is then bounded by the width of one
level rather than by the size of the index, and the pages of each level are read
in increasing page number~--- the same argument as for maintenance, and for the
same reason.

The check runs under share locks and does not block readers or writers. It
therefore observes a moving structure, and its invariants must be ones that hold
of every intermediate state a concurrent modification passes through, not merely
of quiescent states. This is a real restriction on what can be checked, and it is
the reason invariant~1 is stated as coverage rather than as equality: a
concurrent insertion may have already widened a parent key without yet having
inserted the entry that motivated it, and coverage survives that window where
equality would not.

\section{Sibling links and lock coupling}
\label{sec:siblings}

Invariant~4 cannot be checked while holding a lock on one page: the check needs
to see a left neighbour and its right neighbour at once. That means lock
coupling~--- holding one page's lock while acquiring the next~--- which integrity
tools had previously avoided, on the reasonable grounds that it complicates
reasoning about deadlock.

The argument for admitting it here has two halves. Formally, the locks are taken
strictly in one direction, left to right, so no cycle of waits can form with any
other party that also respects that order; and they are share locks, so readers
are unaffected. Practically, sibling verification detects a substantial share of
the corruption actually encountered in the field at very low cost, which is what
justifies the additional complexity. We state both halves because the first alone
would not have been enough to change the practice, and the second alone would not
have been enough to make it safe.

Implemented in \code{39132b784ae}, released in PostgreSQL~13; discussed
in~\cite{thread-rightlink-lockcoupling}.

\section{Normalizing tuples before comparison}
\label{sec:normalize}

The principal way to check that an index agrees with its table is to take an
index tuple and compare it against a tuple built afresh from the table value.
Byte comparison does not work. The same value of a variable-length type may be
represented in several ways~--- compressed or not, with headers of different
lengths~--- and which representation is present depends on the history of
operations that produced it, not on the value. A checker comparing bytes reports
differences that are not differences.

The fix is to normalize both tuples to a canonical representation before
comparing: \code{b1fe8efdf17} for uncompressed variable-length data, and
\code{ab65dfb0fb2} for the differing header sizes of short variable-length data.

The sequel belongs in this paper rather than in a footnote. The normalization
itself then turned out to use the wrong accessor for the size of a
variable-length datum (\code{da1eff08a5b}), which is to say that the code written
to remove a source of false positives introduced one. We return to this in
Section~\ref{sec:checker}.

\section{Running the check on a replica}
\label{sec:replica}

Verification reads the whole index, and a system serving users is the wrong place
to do that. The practical requirement is therefore that the check run on a
physical replica, which shifts the load off the primary and makes it affordable
to run often. Two things had to be corrected for that to work.

Relations that are not fully logged have no meaningful content on a replica, and
checking them produces reports that mean nothing; they are skipped
(\code{6754fe65a4c}). And a check that needs a consistent view of the data must
obtain one under recovery conditions, where the usual mechanism does not apply
unchanged (\code{1f28982e409}).

There is a subtler dependency. On a replica the structure is the product of
replaying the write-ahead log, so the correctness of the check is contingent on
the correctness of that replay. Without the corresponding fixes, the check
reports corruption that exists only in the replica's reconstruction of the index.
This is worth stating plainly: verification on a replica checks the primary's
index and the replay path at once, and does not by itself tell you which of them
is at fault.

\section{Making corruption a machine-readable event}

A report of corruption is only useful to an operator if a program can act on it.
Distinguishing corruption by the text of an error message is not something one
can build automation on. Corruption reports were therefore given error codes of
their own (\code{8ec97e78a77}), so that tooling can recognize the condition and,
for instance, take a node out of service without a human reading a log.

This is the smallest change discussed in this paper and, measured by what it
enables, not the least useful: it is what turns verification from a thing a
person runs during an incident into a thing a system runs continuously.

\section{The checker needed checking}
\label{sec:checker}

Verification of the inverted index shipped in PostgreSQL~18
(\code{14ffaece0fb}, with common routines factored out in \code{d70b17636dd}).
Since then four defects have been fixed in it: posting tree checks
(\code{0cf205e122a}), the parent key check (\code{cdd1a431f21}), the ordering of
entry-tree items (\code{0b54b392334}), and a memory leak in the parent key
consistency pass reported from outside~\cite{thread-gin-amcheck-leak}. Support
for the generalized search tree is partial and still in
review~\cite{thread-amcheck-gist-2025}, seven years after the first attempt at
it~\cite{thread-amcheck-gist-2018}.

We report this deliberately, because the shape of it matters, and it is not the
shape we expected before looking.

Three of the four were not false alarms. They were the opposite: checks that
silently did not check. The clearest is the parent key check
(\code{cdd1a431f21}). It guarded its work with a test on the page's right link,
written as a test for the link being zero~--- but an absent right link is
recorded as the largest representable block number, not as zero. The guard was
therefore almost always false, the parent key was left unset, and no comparison
against it ever ran. The function reported success on indexes it had not
examined. The ordering check for the entry tree skipped comparisons across
attributes of a multi-column index (\code{0b54b392334}); the posting tree check
invalidated parent keys for more pages than necessary and so declined to check
them (\code{0cf205e122a}). In each case the tool answered ``no corruption
found'' without having looked.

That is precisely the failure mode of Section~1, reappearing one level up. The
index answers a query with a subset of the correct answer; the checker answers
the question ``is this index sound'' with a subset of the checks it claims to
perform. Both are silent, both are invisible to functional tests~--- a test that
runs the checker on a good index and expects no complaint passes either way~---
and both are found only by someone reading the code or by an independent
implementation disagreeing. It was in fact a reader: the posting tree issues were
reported from outside by Arseniy Mukhin, as was the memory leak in the parent key
pass~\cite{thread-gin-amcheck-leak}. The normalization defect
(\code{da1eff08a5b}), a variable-length size accessor used before the header size
was known, was reported by Andres Freund.

We had expected, before assembling this list, that the dangerous defect in an
integrity tool would be the false positive: an alarm on a sound index sends the
operator hunting for damage that does not exist, and after that happens twice the
tool stops being run, at which point its detection rate is zero however good it
is. That reasoning still seems right to us, and it is not what the record shows
happened. What happened is that the tool was too quiet, not too loud, and being
too quiet is harder to notice by construction: nobody investigates a clean bill
of health.

It follows that a verification tool is subject to the argument it exists to make.
Its own soundness is a structural property that no functional test of its callers
can observe, and the evidence for it is the same kind of evidence: deployment
across many key types and workloads, over time, plus readers. Seven years from
the first patch~\cite{thread-amcheck-gist-2018} to partial support in review is
not obviously too slow for something that has to meet that standard.

\section{Evaluation}

% TODO. Машина сейчас занята другим агентом; измерения отложены.
% План:
% 1. Накладные расходы проверки: время и число прочитанных страниц относительно
%    размера индекса, для gist_index_check и gin_index_check, при холодном и
%    тёплом кэше. Первичная величина — страницы, время вторично (урок P2).
% 2. Расход памяти при обходе по уровням против обхода сверху вниз — прямая
%    проверка утверждения §\ref{sec:parent} об ограничении шириной уровня.
% 3. Доля обнаруживаемых повреждений на наборе искусственно повреждённых
%    индексов: порча ключа родителя, правой ссылки, отметки удаления, порядка
%    элементов. По инварианту на каждый вид порчи — какой из них срабатывает.
% 4. Проверка на реплике против проверки на основной системе: влияние на
%    задержку пользовательской нагрузки.

\section{Limitations}

Verification of the generalized search tree is partial: what is committed covers
the inverted index, and the tree is still in review. The check runs concurrently
and therefore verifies only invariants that hold of every intermediate state, so
properties true merely of quiescent structures are outside its reach. Coverage of
key types rests on the operator class interface being used as intended; an
operator class whose \code{consistent} is itself wrong will be found consistent
with itself. And on a replica the check cannot separate a defect in the index
from a defect in log replay.

\section{Conclusion}

The failures that matter most in an index access method are the ones that do not
look like failures. Every operation succeeds, the structure quietly stops
agreeing with the table, and nothing in the ordinary path of the system is in a
position to notice. Testing behaviour cannot reach this class, because the
behaviour is correct at every step.

What reaches it is stating the invariants of the structure and checking the bytes
against them. For an extensible framework this can be done without knowing what a
key means, using only the interface the search itself uses, which is what makes
the checker apply to key types written after it. Doing it on a running system
constrains which invariants are available~--- only those that hold of every
intermediate state~--- and doing it on a replica makes it affordable enough to
run continuously rather than during incidents.

The last finding is the one we would least like to be true. The checker is not
exempt from its own argument: its soundness is exactly the kind of property it
was built to check, and it failed in exactly the way it was built to detect.
Three of the four defects found in its first year were checks that quietly did
not run and reported success on indexes they had not examined~--- silent
incompleteness, one level up from the silent incompleteness the tool exists to
catch. We had expected the characteristic failure of an integrity tool to be the
false alarm, which destroys trust and gets the tool switched off. The record says
otherwise, and says something less comfortable: the tool was too quiet, and being
too quiet is the harder fault to find, because nobody investigates a clean bill
of health.
